% ============================================================
% f2014_shanghai_arts_q17_square.tex
% 2014 上海卷（文科）第 17 题：四个小正方形拼成的大正方形
% 来源 PDF：普通高考/2014/2014上海文.pdf
% 原页：第 1 页（A3），印面 917，选择题第 17 题题干下方
% 大正方形边长 2，内含两条中位线；九个格点中 A、B 为大正方形
% 左下、左上顶点（AB 为左边），其余七个格点标为 P1—P7：
%   P1 左边中点，P2 下边中点，P3 中心，P4 上边中点，
%   P5 右下顶点，P6 右边中点，P7 右上顶点。
% ============================================================
\begin{tikzpicture}[>=Stealth, line width=0.65pt, font=\small, scale=2.1,
  line cap=round, line join=round]
  % 大正方形及两条中位线
  \draw (0,0) rectangle (2,2);
  \draw (1,0) -- (1,2);
  \draw (0,1) -- (2,1);

  % 九个格点
  \fill (0,0) circle (0.75pt);
  \fill (0,1) circle (0.75pt);
  \fill (0,2) circle (0.75pt);
  \fill (1,0) circle (0.75pt);
  \fill (1,1) circle (0.75pt);
  \fill (1,2) circle (0.75pt);
  \fill (2,0) circle (0.75pt);
  \fill (2,1) circle (0.75pt);
  \fill (2,2) circle (0.75pt);

  % 顶点标注
  \node[below left=-1pt] at (0,0) {$A$};
  \node[above left=-1pt] at (0,2) {$B$};
  \node[left=-1pt] at (0,1) {$P_1$};
  \node[below=-1pt] at (1,0) {$P_2$};
  \node[above right=-2pt] at (1,1) {$P_3$};
  \node[above=-1pt] at (1,2) {$P_4$};
  \node[below right=-1pt] at (2,0) {$P_5$};
  \node[right=-1pt] at (2,1) {$P_6$};
  \node[above right=-1pt] at (2,2) {$P_7$};
\end{tikzpicture}
